% =============================================================
% chapters/abbreviations.tex
% Перелік умовних скорочень (за вимогою Положення КДПУ §4 для робіт
% з великою кількістю абревіатур, ≥10 одиниць).
% =============================================================

\begin{center}
\textbf{\Large ПЕРЕЛІК УМОВНИХ СКОРОЧЕНЬ}
\end{center}

\vspace{0.5em}

\begin{flushleft}

\noindent\textbf{АІКОМ}~-- автоматизована інформаційна комплексна освітня мережа\\[2pt]
\noindent\textbf{ВПО}~-- внутрішньо переміщені особи\\[2pt]
\noindent\textbf{ВПУ}~-- вищі професійні училища\\[2pt]
\noindent\textbf{ДІСО}~-- державна інформаційна система освіти\\[2pt]
\noindent\textbf{ЄДЕБО}~-- Єдина державна електронна база з питань освіти\\[2pt]
\noindent\textbf{ЗПО}~-- заклад професійної освіти\\[2pt]
\noindent\textbf{ІСУО}~-- інформаційна система управління освітою\\[2pt]
\noindent\textbf{ЛОДА}~-- Луганська обласна державна адміністрація\\[2pt]
\noindent\textbf{МОН}~-- Міністерство освіти і науки України\\[2pt]
\noindent\textbf{МТБ}~-- матеріально-технічна база\\[2pt]
\noindent\textbf{НАПН}~-- Національна академія педагогічних наук України\\[2pt]
\noindent\textbf{НМТ}~-- національний мультипредметний тест\\[2pt]
\noindent\textbf{НМЦ ПТО}~-- Навчально-методичний центр професійно-технічної освіти (історична офіційна назва установи у~Луганській~області)\\[2pt]
\noindent\textbf{ОДА}~-- обласна державна адміністрація\\[2pt]
\noindent\textbf{ПО}~-- професійна освіта\\[2pt]
\noindent\textbf{УЦОЯО}~-- Український центр оцінювання якості освіти\\[2pt]

\vspace{0.5em}

\noindent\textbf{ABAC}~-- attribute-based access control (керування доступом на~основі атрибутів)\\[2pt]
\noindent\textbf{API}~-- application programming interface (програмний інтерфейс додатка)\\[2pt]
\noindent\textbf{CDN}~-- content delivery network (мережа доставки контенту)\\[2pt]
\noindent\textbf{CMS}~-- content management system (система керування контентом)\\[2pt]
\noindent\textbf{CSS}~-- cascading style sheets (каскадні таблиці стилів)\\[2pt]
\noindent\textbf{DSR}~-- design science research (наукове проєктування)\\[2pt]
\noindent\textbf{EQAVET}~-- European Quality Assurance in Vocational Education and Training\\[2pt]
\noindent\textbf{ETF}~-- European Training Foundation\\[2pt]
\noindent\textbf{ETL}~-- extract-transform-load (вилучення-перетворення-завантаження даних)\\[2pt]
\noindent\textbf{EU4Skills}~-- програма ЄС «Кращі навички для сучасної України»\\[2pt]
\noindent\textbf{HTML}~-- hypertext markup language (мова розмітки гіпертексту)\\[2pt]
\noindent\textbf{HTTPS}~-- hypertext transfer protocol secure (захищений протокол передавання гіпертексту)\\[2pt]
\noindent\textbf{IT}~-- information technology (інформаційні технології)\\[2pt]
\noindent\textbf{KPI}~-- key performance indicator (ключовий показник ефективності)\\[2pt]
\noindent\textbf{LMS}~-- learning management system (система керування навчанням)\\[2pt]
\noindent\textbf{SHA-256}~-- secure hash algorithm 256-bit (криптографічна геш-функція)\\[2pt]
\noindent\textbf{SQL}~-- structured query language (мова структурованих запитів)\\[2pt]
\noindent\textbf{SSL}~-- secure sockets layer (рівень захищених сокетів)\\[2pt]
\noindent\textbf{SUS}~-- system usability scale (шкала юзабіліті системи)\\[2pt]
\noindent\textbf{WCAG 2.1 AA}~-- web content accessibility guidelines 2.1 рівень AA (рекомендації щодо доступності веб-контенту)

\end{flushleft}

\clearpage
